\documentclass[11pt,a4paper]{article}
\usepackage[margin=25mm]{geometry}
\usepackage{amsmath,amssymb,amsthm,hyperref}
\newtheorem{theorem}{Theorem}
\newcommand{\gen}[1]{\langle#1\rangle}
\title{Completion of the omission on pairs}
\author{Human Compression project: a component of the L044 recovery}
\date{}
\begin{document}\maketitle
We acquire the complete proofs of \emph{Centralizers of a primitive
action on pairs}, \emph{An alternating group on the pairs},
\emph{Faithful symmetric lifts on pairs}, and \emph{The central
kernel in a symmetric action on pairs}. The last two are separate
parts of the symmetric case; their full hypotheses and finite
premises remain in force. We also acquire the complete finite
maxima through degree fourteen and the retained L047 maximal-block
decomposition. Each component has passed a distinct mathematical review. The
complete proofs, rather than their result statements alone, are
retained dependencies of this assembly.

\begin{theorem}
Let $A_n=\gen{\rho_0,\ldots,\rho_{r-1}}$ have a string
C-marking in its natural action, where $n\ge12$.
Suppose some maximal omission is transitive and has a maximal
block system consisting of pairs. Assume explicitly the alternating
rank bound in every strictly smaller degree $12\le q<n$.
Then $r\le\lfloor(n-1)/2\rfloor$.
\end{theorem}
\begin{proof}
Write $n=2m$. Degrees twelve and fourteen follow from the complete
finite classification. We may therefore suppose $m\ge8$ and,
for a contradiction, $r\ge m$. The block image $P$ is primitive,
because the chosen block system is maximal. Choose an
inclusion-minimal distinguished family $L$ generating $P$,
and let $C$ be the other marks of the omission commuting with
all of $L$, with $R$ the remaining marks.

The complete centralizer component excludes the disconnected
projected diagram: then $m\ge60$ and
$r\le2\log_2m+6<m$. In the connected case it establishes
that $L$ is an interval, $|C|\le1$, and $|R|\le1$.
This last inequality uses the actual pair-centralizer theorem
to make the omitted label a neighbour of $L$; it does not
assume that this omitted permutation acts on the pairs.
If $P\ne A_m,S_m$, the corrected primitive independent-string
bound gives $|L|\le m-4$, hence
$r=1+|L|+|C|+|R|\le m-1$, already a contradiction.

For $P=A_m$, the complete alternating-pair theorem gives the
same bound. Its invocations of alternating rank use either the
exact finite maxima or degrees $12\le q<2m=n$, which are
precisely the standing strong-induction hypotheses.
For $P=S_m$, the faithful component first excludes $|L|=m-1$
and every faithful projection of $\gen L$. It leaves exactly
$m$ even, $r=m$, $|L|=m-2$, $C=\varnothing$, $|R|=1$,
with the central kernel $\gen z$ inside $\gen L$ and full
even-weight kernel $E$ inside the omission. After reversal the
missing and extra marks lie at opposite ends of $L$.
The complete central-kernel component excludes this exact
configuration, including its independent degree-eight boundary.
All primitive block images have therefore been exhausted.
\end{proof}

\paragraph{Attribution and scope.}
This assembles the three block-image cases of source Section~3.2.
The source module reduction, L047 decomposition and acquired
classifications retain their attribution; the actual centralizers,
faithful endpoint arguments and string-only core argument are
identified in the component proofs. No quotient IP or block action
of a block-breaking generator is inferred. The theorem is a branch
of the principal proof, not a whole-original recovery by itself.
\end{document}
